% Study sheet: Volume II, Chapter 3, Network Fabrics.
% Numbers: \val{...} from network_fabrics_studysheet_numbers.tex, written by network_fabrics_studysheet_numbers.py.
\documentclass{studysheet}
\input{network_fabrics_studysheet_numbers}
\sheetvolume{2}
\sheetchapter{3}{Network Fabrics}
\sheetsources{Numbers from the mlsysim registry and the chapter's scenarios; link rates in Gb/s are bits, GB/s bytes; assumed inputs are marked.}

\begin{document}
% ===================================================================== FRONT
\frontside
\begin{bigidea}
At fleet scale the network is the \hl{synchronization backbone}: every step waits at a barrier for its
slowest required transfer. In C\textsuperscript{3} terms the fabric carries communication, and the barrier
turns its delays into coordination. Judge it by \hl{collective completion time and its tail}, not port speed.
\end{bigidea}

\begin{leftcol}
\sheetsec{Key ideas}
\begin{ideas}
\idea{Bulk-synchronous training inverts traffic.}{Few large, synchronized flows meet at a barrier; the
  slowest sets the step, and static hashing can pin two to one link.}
\idea{Two costs per message.}{Fixed startup latency plus per-byte time: control messages are
  latency bound, gradient shards bandwidth bound.}
\idea{The cut caps collectives.}{Bisection bandwidth spans the narrowest cut halving the
  cluster; oversubscription divides it by its ratio.}
\idea{Topology must match traffic.}{Fat-trees give full bisection; rails give same-rank gradient
  traffic its own switches; a torus suits neighbor rings, not all-to-all.}
\idea{Lossless moves the risk.}{Credits (InfiniBand) or PAUSE frames (PFC) prevent drops and ECN
  marks slow senders, but pauses spread to victim flows (a storm) or freeze a cycle of full buffers.}
\idea{Faults hide as silent waste.}{A link that trains at a lower rate or on fewer lanes stays up at
  reduced bandwidth; full rate with rising pause counters points to congestion.}
\end{ideas}

\vspace{10pt}
\sheetsec{Sizing and stress rules}
\begin{sheettable}{@{}L c L@{}}
\thead{Rule} & \thead{Formula} & \thead{Radix-\val{radix} / chapter case}\\ \midrule
Two-tier fat-tree hosts & $\SwitchRadix^2/2$ & \val{twotier}; longest path three switches\\
Three-tier fat-tree hosts & $\SwitchRadix^3/4$ & \val{threetier}; longest path five switches\\
Incast: buffer fills in & $\dfrac{\EgressBufferCapacity}{\EgressArrivalRate-\EgressServiceRate}$ & \val{incsend}\,$\times$\,\val{ndrbytes}~GB/s into one \val{ndrbytes}~GB/s port, \val{incbuf} MB: \val{incfill}\,$\MicroUnitPrefix$s, under one \val{incrtt}\,$\MicroUnitPrefix$s feedback round trip\\
\end{sheettable}
\vspace{3pt}
{\fontsize{8.4}{10.4}\selectfont\color{muted}Radix $\SwitchRadix$: ports per switch. At one port:
arrivals $\EgressArrivalRate$, drain $\EgressServiceRate$, buffer $\EgressBufferCapacity$. Separate traffic
classes (virtual lanes, QoS) keep checkpoint bursts out of a collective's buffers.\par}
\end{leftcol}%
\begin{rightcol}
\fcard{Alpha-beta model}%
  {\PointToPointTransferTime(\PointToPointMessageSize) = \MessageStartupLatency + \frac{\PointToPointMessageSize}{\LinkBandwidth}
   \qquad \LatencyBandwidthCrossoverSize = \MessageStartupLatency\LinkBandwidth}%
  {A message of $\PointToPointMessageSize$ bytes pays startup $\MessageStartupLatency$ plus
   $\PointToPointMessageSize/\LinkBandwidth$. Far below $\LatencyBandwidthCrossoverSize$ more bandwidth barely helps;
   far above it, lower latency barely helps.}%
  {Modeled NDR path: $\MessageStartupLatency = \val{alpha}\,\MicroUnitPrefix$s,
   $\LinkBandwidth = \val{ndrbytes}$~GB/s (\val{ndrgbit}~Gb/s), so $\LatencyBandwidthCrossoverSize \approx \val{crossover}$~KB.}%
  {\begin{tikzpicture}[sheet]
     \draw[dim] (24,0) -- (220,0);
     \draw[dim] (24,0) -- (24,90);
     \foreach \d/\tl in {2/100 B,4/10 KB,6/1 MB,8/100 MB}{
       \pgfmathsetmacro\x{24+(\d-2)*28}
       \draw[dim] (\x,0) -- (\x,-2.5);
       \node[small, anchor=north] at (\x,-4) {\tl};}
     \foreach \d/\tl in {0/1\,µs,2/100\,µs,4/10\,ms}{
       \pgfmathsetmacro\y{\d*20}
       \draw[dim] (24,\y) -- (21.5,\y);
       \node[small, anchor=east] at (20,\y) {\tl};}
     \node[small, anchor=north] at (122,-13) {message size (log)};
     \node[small, anchor=south west] at (2,92) {time (log)};
     \draw[gridline, dashed] (24,\val{yalpha}) -- (150,\val{yalpha});
     \fill[accentsoft] (24,0.5) rectangle (\val{xknee},88);
     \edef\abplot{\noexpand\draw[accent, line width=1pt] plot coordinates {\val{abcurve}};}\abplot
     \draw[accent, dashed, line width=0.6pt] (\val{xknee},0) -- (\val{xknee},\val{yknee});
     \filldraw[fill=white, draw=accent, line width=0.6pt] (\val{xknee},\val{yknee}) circle (1.8pt);
     \node[note, anchor=south east] at (\val{xknee}-1,\val{yknee}+12) {$\LatencyBandwidthCrossoverSize$ = \val{crossover} KB};
     \node[small, anchor=north west] at (27,86) {latency bound};
     \node[small, anchor=north west] at (\val{xknee}+3,86) {bandwidth bound};
     \filldraw[fill=accent, draw=accent] (\val{xctrl},\val{yctrl}) circle (2.2pt);
     \node[note, anchor=south] at (\val{xctrl},\val{yctrl}+6) {control \val{ctrlkb} KB};
     \filldraw[fill=accent, draw=accent] (\val{xshard},\val{yshard}) circle (2.2pt);
     \node[note, anchor=east] at (\val{xshard}-5,\val{yshard}+2) {shard \val{shardmb} MB};
   \end{tikzpicture}}
\vspace{4pt}
\fcard{Capacity across the cut}%
  {\BisectionBandwidth = \frac{\DeviceCount}{2}\,\LinkBandwidth
   \qquad \text{relative throughput} = \frac{1}{(1-\CommunicationTimeFraction) + \CommunicationTimeFraction\,\CommunicationSlowdownFactor}}%
  {Per direction, $\DeviceCount$ devices, non-blocking; an oversubscribed tier divides it by its ratio.
   $\CommunicationTimeFraction$: share of step time spent crossing the cut; $\CommunicationSlowdownFactor$: how much
   that time stretches (the ratio, when the cut is saturated). Loss is one minus relative throughput.}%
  {Chapter: \val{cutgpus} GPUs at \val{ndrgbit}~Gb/s: \val{cutfull} TB/s; a \val{cutratio}:1 spine leaves
   \val{cutover}~TB/s. At $\CommunicationTimeFraction$ = \val{twocomm}\%, a \val{tworatio}:1 spine costs \val{twoloss}\% of throughput.}%
  {}
\end{rightcol}

% ===================================================================== BACK
\backside
\begin{answers}
\ans{Finish}{$\LatencyBandwidthCrossoverSize$ = \val{roceal}\,$\MicroUnitPrefix$s $\times$ \val{ndrbytes} GB/s
  = \val{rocecross} KB. Handoff: \val{roceal} + \val{handxfer} = \val{handroce}\,$\MicroUnitPrefix$s, against
  \val{handndr} on NDR (\val{handratio}$\times$ slower). Shard: only \val{sharddiff}\,$\MicroUnitPrefix$s more
  (\val{sharddiffpct}\%): \val{shardtime} ms on both. Startup matters only for small messages (alpha-beta card).}
\ans{1}{\val{qoneside} $\times$ \val{ndrbytes} GB/s = \val{qonefull} TB/s; the spine divides it by \val{qoneratio}: \val{qoneover} TB/s.
  Throughput keeps one over (\val{qonekeep} + \val{qonestretch}) $=$ \val{qonerel}: \val{qoneloss}\% lost (cut card).}
\ans{2}{Two tiers hold $\SwitchRadix$ squared over two hosts: \val{qtwosmallhosts} with radix \val{radix}, too few;
  \val{qtwobighosts} with radix \val{qtwobig}, enough. Radix \val{radix} needs three tiers: more switches,
  five per path instead of three (sizing table).}
\columnbreak
\ans{3}{Arrivals \val{qthreeload} TB/s, drain \val{ndrbytes} GB/s: \val{qthreebuf} MB $\div$ \val{qthreenet} TB/s
  = \val{qthreefill}\,$\MicroUnitPrefix$s, about \val{qthreertts} round trips. Marks on the first queued bytes
  can act in time. A separate class for checkpoint traffic protects the collective (sizing table).}
\ans{4}{Read the pair's link state. A rate of \val{qfourhdr} GB/s (HDR) or half its lanes is a degraded
  link, even without symbol errors; full rate with rising pause counters is congestion (key idea on silent waste).}
\end{answers}

\begin{leftcol}
\sheetsec{Worked example}
\begin{practicebox}\boxtitle{Will a faster link speed up this step's messages?}
A pipeline \val{ctrlkb} KB control message and a \val{shardmb} MB gradient shard on the modeled NDR path
($\MessageStartupLatency = \val{alpha}\,\MicroUnitPrefix$s, $\LinkBandwidth = \val{ndrbytes}$~GB/s).
\begin{steps}
\item Crossover: \val{alpha}\,$\MicroUnitPrefix$s $\times$ \val{ndrbytes} GB/s $=$ \res{\val{crossover} KB}.
  The control message is below it, the shard far above.
\item Control: \val{alpha} + \val{ctrlkb} KB$/$\val{ndrbytes} GB/s $=$ \val{alpha} + \val{ctrlxfer} $=$
  \res{\val{ctrltime}\,$\MicroUnitPrefix$s}; bytes are \val{ctrlshare}\% of it.
\item Shard: \val{shardmb} MB$/$\val{ndrbytes} GB/s $=$ \res{\val{shardtime} ms}; startup is negligible.
\item Double $\LinkBandwidth$ to \val{betadouble} GB/s: control \val{ctrltimetwo}\,$\MicroUnitPrefix$s
  (\val{ctrlsaving}\% less time); shard \val{shardtimetwo} ms (\val{shardspeedup}$\times$ faster).
\end{steps}
\takeaway{Compare each message with $\LatencyBandwidthCrossoverSize$ before buying bandwidth.}
\end{practicebox}

\vspace{12pt}
\sheetsec{Your turn: finish it}
\begin{practicebox}\boxtitle{Does the transport matter at the same line rate?}
\assumed{} a RoCE path at the same \val{ndrgbit}~Gb/s ($\LinkBandwidth = \val{ndrbytes}$~GB/s) with
$\MessageStartupLatency = \val{roceal}\,\MicroUnitPrefix$s; a \val{handkb} KB pipeline handoff and the \val{shardmb} MB shard.\par\vspace{4pt}
\textbf{Given:} $\LatencyBandwidthCrossoverSize = \val{roceal}\,\MicroUnitPrefix\text{s} \times \val{ndrbytes}$~GB/s.\par
\textbf{You:} finish the crossover, time both messages on RoCE and on NDR, and say which messages the
transport affects.
\end{practicebox}
\end{leftcol}%
\begin{rightcol}
\sheetsec{Common mistakes}
\mistake{``Going from HDR (\val{hdrgbit}~Gb/s) to NDR (\val{ndrgbit}~Gb/s) doubles the rate, so every message gets twice as fast.''}%
  {A \val{fallkb} KB message drops from \val{fallhdr} to \val{fallndr}\,$\MicroUnitPrefix$s, only \val{fallspeed}$\times$
   faster (\val{fallless}\% less time) with the same \val{alpha}\,$\MicroUnitPrefix$s startup, which dominates.}{}
\mistake{``PFC makes Ethernet lossless, so congestion is solved.''}%
  {It turns loss into congestion spreading: PAUSE frames climb tier by tier and stall victim flows.
   A storm ends when the slow receiver recovers; a cyclic deadlock does not.}{}
\mistake{``Adaptive routing will fix our oversubscribed spine.''}%
  {It moves flows off colliding paths but adds no capacity; a full cut needs placement that avoids it, or more capacity.}{}

\vspace{2pt}
\sheetsec{Check yourself}
\begin{qs}
\item \assumed{} a \val{qonegpus}-GPU non-blocking fat-tree with \val{ndrgbit}~Gb/s links. Its bisection
  bandwidth? What share of throughput would a \val{qoneratio}:1 spine cost a job spending \val{qonecomm}\% of
  each step on cut traffic?
\item \assumed{} \val{qtwogpus} GPUs, one port each, on a two-tier non-blocking fat-tree. Which rule tells you whether
  radix-\val{radix} or radix-\val{qtwobig} switches suffice, and what does the smaller radix cost?
\item \assumed{} \val{qthreesend} GPUs write a checkpoint at \val{ndrbytes} GB/s each into one \val{ndrgbit}~Gb/s
  port (\val{qthreebuf} MB buffer) that a training collective also uses. How soon does it fill? Can
  \val{incrtt}\,$\MicroUnitPrefix$s feedback act in time? What protects the collective?
\item Step time drifts up a few percent with no errors; one node pair tests at
  \val{qfourmeas} GB/s instead of \val{ndrbytes}. What evidence separates a degraded link from congestion?
\end{qs}
\end{rightcol}
\end{document}
